\section{Appendix: Remote GPU Access (VirtualGL + TurboVNC)}
\label{sec:appendix_remote_gpu}

This section describes how to set up remote GPU-accelerated OpenGL rendering on a Linux server (demonstrated on AlmaLinux~9), enabling COMPASS to be run remotely with full Geographic View support over a VNC session.

\subsection{Overview}

The setup uses:

\begin{itemize}
\item \textbf{VirtualGL}: Redirects OpenGL calls to the server's physical GPU.
\item \textbf{TurboVNC}: High-speed VNC server that integrates with VirtualGL.
\item \textbf{Desktop environment}: e.g. XFCE or KDE (XFCE is recommended as it is lightweight).
\end{itemize}

The instructions below use AlmaLinux~9 / RHEL~9 commands (\texttt{dnf}). For Debian/Ubuntu-based systems, substitute with \texttt{apt} and the corresponding package names.

\subsection{Server Installation}

\subsubsection{Install EPEL and Required Tools}

\begin{lstlisting}
sudo dnf install epel-release
sudo dnf groupinstall "Development Tools"
\end{lstlisting}

\subsubsection{Install a Desktop Environment}

For XFCE (lightweight, recommended for VNC):

\begin{lstlisting}
sudo dnf groupinstall "Xfce"
\end{lstlisting}

For KDE:

\begin{lstlisting}
sudo dnf groupinstall "KDE Plasma Workspaces"
\end{lstlisting}

Optionally set the default session:

\begin{lstlisting}
echo "exec startxfce4" > ~/.xinitrc
\end{lstlisting}

\subsubsection{Install NVIDIA Drivers}

Install the NVIDIA driver from the CUDA repository:

\begin{lstlisting}
sudo dnf install \
  https://developer.download.nvidia.com/compute/cuda/repos/rhel9/x86_64/cuda-rhel9.repo
sudo dnf install nvidia-driver nvidia-settings nvidia-xconfig
sudo nvidia-xconfig
sudo reboot
\end{lstlisting}

After rebooting, verify the GPU is detected:

\begin{lstlisting}
nvidia-smi
\end{lstlisting}

\subsubsection{Install TurboVNC and VirtualGL}

Download the \texttt{.rpm} packages for TurboVNC and VirtualGL from \url{https://www.turbovnc.org}, then install:

\begin{lstlisting}
sudo dnf install ./turbovnc-*.rpm
sudo dnf install ./virtualgl-*.rpm
\end{lstlisting}

\subsubsection{Configure TurboVNC}

Set the VNC password:

\begin{lstlisting}
vncpasswd
\end{lstlisting}

Start the VNC server to verify it works, then stop it:

\begin{lstlisting}
vncserver :1
vncserver -kill :1
\end{lstlisting}

\subsubsection{Configure VirtualGL}

Run the VirtualGL configuration utility:

\begin{lstlisting}
sudo /opt/VirtualGL/bin/vglserver_config
\end{lstlisting}

Recommended answers:

\begin{itemize}
\item \textbf{Restrict to vglusers group?} Yes
\item \textbf{Automatically add users to vglusers group?} Yes
\item \textbf{Disable fake X server?} No (unless using a real X server for display)
\item \textbf{Display manager restart?} No (manual setup preferred)
\end{itemize}

Then add your user to the \texttt{vglusers} group and re-login:

\begin{lstlisting}
sudo usermod -aG vglusers $USER
\end{lstlisting}

\subsection{Running COMPASS Remotely}

\subsubsection{Start a VNC Session on the Server}

\begin{lstlisting}
vncserver :1
\end{lstlisting}

This starts a VNC session on display \texttt{:1} (port 5901).

\subsubsection{Connect from the Client}

From your local machine, use a VNC viewer (e.g. TigerVNC, RealVNC, or the TurboVNC viewer) to connect to:

\begin{lstlisting}
your-server-ip:1
\end{lstlisting}

\subsubsection{Run COMPASS with VirtualGL}

Once inside the VNC session, open a terminal and run the COMPASS AppImage with \texttt{vglrun}:

\begin{lstlisting}
vglrun ./COMPASS-x86_64_release.AppImage
\end{lstlisting}

Alternatively, when connecting via SSH with VirtualGL forwarding:

\begin{lstlisting}
vglconnect user@your-server
vglrun ./COMPASS-x86_64_release.AppImage
\end{lstlisting}

This redirects all OpenGL rendering to the server's physical GPU, streaming the rendered frames to the VNC client with full OpenGL support.

\subsection{Reference}

\begin{itemize}
\item VirtualGL installation guide: \url{https://virtualgl.org/Documentation/Installation}
\item TurboVNC user guide: \url{https://www.turbovnc.org/Documentation/Guide}
\end{itemize}
